% Avsnitt 2.1-2.6, ur lectures/L02/appendix/a_larger_networks.md (A.1-A.6).

\section{Karnaughdiagram: från sanningstabell till minimerat grindnät}
\label{c2:sec:kmap}

Att läsa av en summa-av-produkter-ekvation direkt ur en sanningstabell, så som
\secref{c1:sec:boolean} gjorde, fungerar alltid men ger sällan det minsta nätet: varje
\code{1}-rad blir en egen AND-term, även när flera rader delar de flesta av sina ingångar och
skulle kunna slås ihop.

Ett \textbf{Karnaughdiagram} (K-diagram) gör de gemensamma mönstren synliga för hand, för
funktioner av 2 till 4 variabler, utan någon formell manipulation av boolesk algebra. Ett
genomarbetat exempel följer; mer övning finns i \secref{c2:sec:exercises}, och klassisk
minimeringsteori (don't-cares, Quine\--McCluskey) ligger utanför bokens ram.

Idén är att lägga ut sanningstabellen som ett rutnät i stället för som en lista, och ordna
rubrikerna så att varje ruta skiljer sig från var och en av sina grannar i sidled och i höjdled i
exakt ett ingångsvärde, även runt kanterna vänster/höger och upp/ner. Den ordningen är
\textbf{Graykod}: för två bitar \code{00, 01, 11, 10}, där bara en bit ändras i varje steg,
inklusive från den sista posten tillbaka till den första.

Sedan:
\begin{enumerate}
  \item Skriv en \code{1} i varje ruta där utgången är \code{1}, och lämna resten tomma.
  \item Rita rektangulära grupper runt \code{1}:orna:
  \begin{itemize}
    \item Varje grupp måste vara en rektangel vars bredd och höjd var för sig är en tvåpotens
      (1, 2 eller 4), så att också dess storlek är en tvåpotens.
    \begin{itemize}
      \item en grupp om 4 är alltså antingen en 1x4-linje eller ett 2x2-block, och båda är
        tillåtna.
      \item rutor som bara möts i ett hörn är ingen rektangel, så de kan inte grupperas.
    \end{itemize}
    \item Grupper får överlappa, och får gå runt en kant av rutnätet.
  \end{itemize}
  \item Gör varje grupp så stor som möjligt, och täck varje \code{1} minst en gång. En grupp om 1
    behåller alla sina ingångar i termen, en grupp om 2 stryker en, en grupp om 4 stryker två, och
    så vidare.
  \begin{itemize}
    \item Använd sedan så få grupper som räcker: stryk varje grupp vars rutor allihop redan täcks
      av en annan. Ingenting ovan förbjuder en överflödig grupp, och en överflödig grupp är en
      överflödig AND-term i svaret, korrekt men inte minimal.
  \end{itemize}
  \item Behåll för varje grupp bara den eller de ingångar som har samma värde i varje ruta i
    gruppen, stryk resten, och skriv det som en AND-term:
  \begin{itemize}
    \item skriv en behållen ingång direkt om det konstanta värdet är \code{1}, och primmad om det
      är \code{0}.
    \item en grupp som bara täcker rutor där $B = 0$ bidrar alltså med termen $B'$.
  \end{itemize}
  \item $X$ är alla gruppernas termer OR:ade med varandra.
\end{enumerate}

\subsection*{Genomarbetat exempel}
Tre ingångar \code{A}, \code{B}, \code{C}, en utgång \code{X}:

\begin{narrowtable}{cccc}
  \thead{A} & \thead{B} & \thead{C} & \thead{X} \\
  \midrule
  0 & 0 & 0 & 0 \\
  0 & 0 & 1 & 1 \\
  0 & 1 & 0 & 0 \\
  0 & 1 & 1 & 1 \\
  1 & 0 & 0 & 0 \\
  1 & 0 & 1 & 1 \\
  1 & 1 & 0 & 1 \\
  1 & 1 & 1 & 1 \\
\end{narrowtable}

Lägg ut den med \code{AB} i Graykod nedåt i raderna och \code{C} i sidled i kolumnerna, och fyll i
en \code{1} överallt där $X = 1$:

\lecfig[0.42]{lectures/L02/appendix/images/karnaugh_filled.png}{Karnaughdiagram för \code{X},
ifyllt från sanningstabellen.}{c2:fig:kmap1}

Hela $C = 1$-kolumnen är fyra \code{1}:or, en grupp om 4. Varje ruta delar $C = 1$ och inget annat
är konstant ($A$ och $B$ varierar båda), så den reduceras till envariabeltermen $C$:

\lecfig[0.42]{lectures/L02/appendix/images/karnaugh_group_c.png}{Samma diagram med den fyra rutor
stora $C = 1$-gruppen markerad.}{c2:fig:kmap2}

En \code{1} är fortfarande inte täckt, på raden $AB = 11$. Gruppera den med sin granne till höger,
som $C$-gruppen redan täcker. Båda rutorna i paret delar $A = 1$ och $B = 1$, och $C$ varierar, så
det reduceras till $AB$:

\lecfig[0.42]{lectures/L02/appendix/images/karnaugh_groups.png}{Samma diagram med den två rutor
stora $AB = 11$-gruppen tillagd.}{c2:fig:kmap3}

Varje \code{1} är nu täckt; rutan $AB = 11$, $C = 1$ ligger i båda grupperna, och överlapp är i
sin ordning. Summeringen ger:

\[ X = AB + C \]

Att läsa sanningstabellen direkt, på \chapref{c1:ch}:s vis, hade gett fem AND-termer, en per
\code{1}-rad. K-diagrammet hittade en likvärdig ekvation med två termer, vilket är hela värdet av
att göra det här för hand.

\subsection*{Att realisera nätet}
$X = AB + C$ är två grindar: en \code{AND} som beräknar $AB$ och matar en \code{OR} vars andra
ingång är $C$.

\lecfig[0.55]{lectures/L02/appendix/images/net.png}{Grindnät för det minimerade uttrycket.}
{c2:fig:net}

Du bygger det här nätet för hand i CircuitVerse under passet, precis som du byggde
\code{or_gate}s nät i \chapref{c1:ch}, och översätter det till VHDL i \secref{c2:sec:tovhdl}.

% ------------------------------------------------------------------------------------------
\section{Att översätta Karnaughdiagrammets nät till VHDL}
\label{c2:sec:tovhdl}

$X = AB + C$ översätts rakt av till VHDL. Tre ingångar, en utgång:

\begin{vhdlcode}
library ieee;
use ieee.std_logic_1164.all;

entity gate_network is
    port(a, b, c: in std_logic;
         x      : out std_logic);
end entity;
\end{vhdlcode}

Arkitekturen använder en intern signal \code{ab} för att hålla \code{AND}-grindens utgång, och
speglar tvågrindsnätet ett till ett:

\begin{vhdlcode}
architecture behaviour of gate_network is
signal ab: std_logic;
begin
    -- AND gate: AB.
    ab <= a and b;

    -- OR gate: AB + C.
    x <= ab or c;
end architecture;
\end{vhdlcode}

Utan mellansignalen slås det ihop till \code{x <= (a and b) or c;}. Båda är lika giltiga och
syntetiseras till samma hårdvara; använd den som är mest läsbar för nätet du har framför dig.

Den byggbara versionen använder den hopslagna formen, eftersom det vid två grindar inte finns
något att vinna på att namnge mellansignalen:

\vhdlfile{lectures/L02/gate_network/gate_network.vhd}

Öppna den vid sidan av versionen ovan och övertyga dig om att de beskriver samma krets; det är
påståendet det här avsnittet gör, och det kostar inget att kontrollera.

För att köra den på hårdvara, tilldela portarna till strömbrytare och en lysdiod med DE0-CV:s pin
planner (bilaga~\ref{app:quartus}). Ingenting här beror på specifika pinnummer.

% ------------------------------------------------------------------------------------------
\section{\code{std_logic_vector}: mer än en ledning}
\label{c2:sec:vectors}

Allt hittills har varit \code{std_logic}, en ledning som bär en bit. Att bunta ihop flera ledningar
under ett namn är vad \code{std_logic_vector} är till för.

\begin{vhdlcode}
signal count: std_logic_vector(3 downto 0);
\end{vhdlcode}

Det deklarerar fyra ledningar: \code{count(3)} ner till \code{count(0)}. Intervallet är
\code{3 downto 0} snarare än \code{0 to 3} så att biten längst till vänster är den mest
signifikanta, så som du skulle skriva talet på papper. Boken använder \code{downto} överallt, och
det gör nästan all VHDL du kommer att möta.

En vektor är en porttyp som vilken annan som helst:

\begin{vhdlcode}
entity display is
    port(number: in  std_logic_vector(3 downto 0);
         hex   : out std_logic_vector(6 downto 0));
end entity;
\end{vhdlcode}

\subsection*{Värden}
En enskild bit tar enkla citattecken, en vektor dubbla:

\begin{vhdlcode}
x     <= '1';                -- std_logic: one bit
count <= "1010";             -- std_logic_vector(3 downto 0): four bits
count <= (others => '0');    -- every bit '0', whatever the width
\end{vhdlcode}

Den sista är ett \emph{aggregat}: \code{(others => '0')} betyder ”alla återstående bitar är
\code{'0'}”, och eftersom ingen namngavs individuellt är det allihop. Det fortsätter fungera när
bredden ändras, och det dyker upp genom hela resten av boken.

\subsection*{Indexering och slicing}
En bit ur en vektor är en \code{std_logic}; en sammanhängande följd är en \emph{slice}, som själv
är en vektor:

\begin{vhdlcode}
carry     <= count(3);
high_half <= input(7 downto 4);
\end{vhdlcode}

Slicing är hur en bred port matar flera smalare, vilket är vad \secref{c2:sec:submodules}:s
tvåsiffriga display gör med sina åtta strömbrytaringångar.

En detalj är inte vad en mjukvaruutvecklare väntar sig. \code{input(7 downto 4)} har
indexintervallet \code{7 downto 4}, inte \code{3 downto 0}. Kopplad till en port deklarerad
\code{std_logic_vector(3 downto 0)} matchas de två \textbf{positionellt, från vänster till höger},
inte efter indexnummer: \code{input(7)} driver bit 3, \code{input(6)} driver bit 2, och så vidare.
Längderna måste stämma överens. Numreringen behöver inte det.

\subsection*{En vektor är ett knippe bitar, inte ett tal}
\code{count <= "1010"} gör inte \code{count} lika med tio. \code{std_logic_vector} säger ingenting
om vad bitarna \emph{betyder}, så den erbjuder bara de operationer som är meningsfulla för ett
knippe ledningar: de logiska operatorerna tillämpade bit för bit (\code{a and b}, \code{not a},
\code{a xor b}), jämförelse (\code{count = "1010"}) och konkatenering (\code{"10" & "11"} är
\code{"1011"}).

Aritmetik står inte på den listan. \code{count + 1} går inte att kompilera, eftersom ingenting har
sagt om de fyra ledningarna är ett teckenlöst tal, ett tecknat, eller fyra orelaterade signaler.

Boken går inte runt det med en numerisk vektortyp. Allt som räknas är ett vanligt heltal med ett
intervall:

\begin{vhdlcode}
signal count: natural range 0 to 9;
\end{vhdlcode}

som du kan addera till och jämföra direkt, och som syntesen gör om till exakt de vippor
intervallet behöver. Vektorer förblir vad de är: knippen av ledningar, för saker som verkligen är
knippen av ledningar, som åtta strömbrytare eller sju displaysegment.

De två världarna möts bara vid en gräns, och konverteringarna dyker upp där de behövs:

\begin{vhdlcode}
count_v <= std_logic_vector(to_unsigned(count, 4));   -- integer to vector
count   <= to_integer(unsigned(count_v));             -- vector to integer
\end{vhdlcode}

Båda kommer från \code{ieee.numeric_std}, vilket är varför vissa filer lägger till en tredje
\code{use}-sats.

% ------------------------------------------------------------------------------------------
\section{Multiplexrar}
\label{c2:sec:mux}

En \textbf{multiplexer} (”mux”) kopplar exakt en av flera dataingångar till en enda utgång, vald
av en separat uppsättning \textbf{väljaringångar}. En \code{N}-till-1-mux har \code{N}
dataingångar och $\lceil \log_2 N \rceil$ väljarbitar.

Det enklaste fallet, en 2-till-1-mux med ingångarna \code{A}, \code{B}, väljaren \code{S} och
utgången \code{X}:

\begin{narrowtable}{cc}
  \thead{S} & \thead{X} \\
  \midrule
  0 & B \\
  1 & A \\
\end{narrowtable}

Som boolesk ekvation, $X = S'B + SA$: för varje kombination av ingångar är exakt en AND-term
aktiv, och den termens dataingång släpps igenom.

Större muxar är samma mönster med fler väljarbitar. En 8-till-1-mux med dataingångarna
\code{A}--\code{H} och väljaren \code{S[2:0]}:

\begin{narrowtable}{rc}
  \thead{S[2:0]} & \thead{X} \\
  \midrule
  000 & A \\
  001 & B \\
  010 & C \\
  011 & D \\
  100 & E \\
  101 & F \\
  110 & G \\
  111 & H \\
\end{narrowtable}

Läst som en summa av produkter med en term per rad bidrar raden \code{000} med $A\,S_2'S_1'S_0'$
och resten följer samma form, där var och en grindar sin dataingång med den väljarkombination som
väljer den:

\[ X = A\,S_2'S_1'S_0' + B\,S_2'S_1'S_0 + \dots + H\,S_2S_1S_0 \]

Värt att skriva ut en gång för hand, inte värt att memorera. Mönstret är poängen, och det är
detsamma i vilken storlek som helst.

Som grindar: en \code{AND} med fyra ingångar per dataingång (databiten plus alla tre
väljarbitarna), med väljarbitarna inverterade efter behov, och alla åtta \code{AND}-utgångar
summerade av ett träd av \code{OR}-grindar.

\lecfig[0.72]{lectures/L02/appendix/images/mux_8to1.png}{8-till-1-multiplexer som åtta AND-grindar
som matar ett OR-träd.}{c2:fig:mux8}

Det här är en verklig byggsten. ATmega328P (chippet på en Arduino Uno och Nano) använder en
internt så att dess analogkapabla pinnar, plus en intern bandgapsreferens och en temperatursensor,
kan dela på en enda analog-till-digital-omvandlare, vald av de fyra \code{MUX}-bitarna i dess
\code{ADMUX}-register.

% ------------------------------------------------------------------------------------------
\section{Multiplexrar i VHDL: \code{process} och \code{case}}
\label{c2:sec:process}

Varje sats hittills har varit \textbf{konkurrent}: en stående beskrivning av en ledning, alla
aktiva samtidigt. En multiplexer är den första kretsen som är otymplig att skriva så, och den
behöver två nya konstruktioner:
\begin{itemize}
  \item en \textbf{\code{process}}, ett block vars innehåll körs sekventiellt, uppifrån och ned,
    som vanlig kod. Den körs om varje gång någon signal i dess \textbf{känslighetslista} ändras,
    och ur arkitekturens synvinkel är den fortfarande bara ännu en konkurrent sats. Varje process
    i boken har den här formen; \chapref{c3:ch} återanvänder den för klockad logik, med klockan i
    känslighetslistan i stället för de ingångar som avkodas.
  \item en \textbf{\code{case}-sats}, VHDL:s \code{switch}. Precis som \code{if} och \code{for} är
    den en \emph{sekventiell} sats, tillåten bara inuti en process, och varje värde väljaren kan
    anta måste täckas, vilket \code{when others} garanterar.
\end{itemize}

”Välj en av flera ingångar utifrån en väljare” \emph{är} en \code{case}-sats, nästan ord för ord.
2-till-1-muxen från \secref{c2:sec:mux}:

\begin{vhdlcode}
MUX_PROCESS: process(a, b, sel) is
begin
    case (sel) is
        when '0'    => x <= b;
        when '1'    => x <= a;
        when others => x <= '0';
    end case;
end process;
\end{vhdlcode}

Två saker att notera:
\begin{itemize}
  \item Enbitsvärden tar enkla citattecken, \code{'0'}/\code{'1'}, till skillnad från de dubbla
    citattecken som används för vektorer som \code{"000"}.
  \item \code{when others} är obligatoriskt snarare än defensivt. VHDL kräver att ett \code{case}
    täcker varje värde av väljarens typ, och \code{std_logic} har nio (\secref{c1:sec:vhdl}), så
    \code{'0'} och \code{'1'} ensamma gör aldrig ett \code{case} uttömmande. Det är en språkregel,
    inte ett påstående om ledningen: ett syntetiserat nät bär en \code{0} eller en \code{1} och
    inget annat, och de andra sju värdena är simulatorns, där en signal som ligger kvar på
    \code{'U'} eller \code{'X'} är en bugg snarare än en nivå.
\end{itemize}

\subsection*{Genomarbetat exempel: 8-till-1-muxen}
Samma form med en trebitars väljare och åtta grenar, en per rad i \secref{c2:sec:mux}:s
sanningstabell. Dataingångarna \code{A}--\code{H} blir vektorn \code{inputs(7 downto 0)}, med
\code{A} som \code{inputs(7)} ner till \code{H} som \code{inputs(0)}, så väljaren \code{"000"}
kopplar fram \code{inputs(7)} och \code{"111"} kopplar fram \code{inputs(0)}.

\vhdlfile{lectures/L02/mux_8to1/mux_8to1.vhd}

Jämför den med \secref{c2:sec:mux}:s ekvation med åtta termer för samma krets. Ekvationen
beskriver grindarna; \code{case}-satsen beskriver \emph{avsikten} och låter syntesverktyget ta
fram grindarna. Båda är korrekta och båda syntetiseras till samma hårdvara. Det glappet, mellan
att beskriva struktur och att beskriva beteende, är det mesta av vad ett hårdvarubeskrivande språk
ger dig.

% ------------------------------------------------------------------------------------------
\section{Att bygga en design av submoduler}
\label{c2:sec:submodules}

Allt hittills har varit en entitet med en arkitektur. Det slutar räcka i samma stund en design
behöver samma block två gånger.

DE0-CV har sex 7-segmentsdisplayer, och du vill driva två av dem, var och en visande en
hexadecimal siffra. Logiken som gör ett fyrabitars tal till sju segmentsignaler är densamma för
båda. Du skulle kunna skriva den två gånger inuti en arkitektur och ingenting skulle hindra dig,
men den andra kopian är en belastning: varje rättning måste göras i båda, och den dag de glider
isär är den dag du förlorar en kväll på att lista ut vilken av dem som är fel.

Skriv blocket en gång, som en egen entitet, och \emph{instansiera} det två gånger.

\subsection*{Submodulen}
En vanlig modul. Ingenting hos den säger ”jag är en submodul”:

\begin{vhdlcode}
entity display is
    port(number: in  std_logic_vector(3 downto 0);
         hex   : out std_logic_vector(6 downto 0));
end entity;
\end{vhdlcode}

Om den slutar som en hel design eller som en del av en större avgörs av den som instansierar den,
inte av modulen.

\subsection*{Att instansiera den}
Med \code{entity work.<name>}, inuti en annan arkitektur:

\begin{vhdlcode}
entity hex_display is
    port(input     : in  std_logic_vector(7 downto 0);
         hex1, hex0: out std_logic_vector(6 downto 0));
end entity;

architecture structure of hex_display is
begin
    display1: entity work.display
        port map(input(7 downto 4), hex1);

    display0: entity work.display
        port map(input(3 downto 0), hex0);
end architecture;
\end{vhdlcode}

\begin{itemize}
  \item \code{display1} och \code{display0} är \textbf{instansetiketter}, som namnger de två
    kopiorna. De måste vara unika inom arkitekturen, och de är inte modulens namn: båda
    instanserna är \code{display}.
  \item \code{entity work.display} väljer modulen. \code{work} är det bibliotek dina analyserade
    filer hamnar i som förval, i både GHDL och Quartus.
  \item \code{port map(...)} kopplar instansens portar till signaler i den omslutande
    arkitekturen, \textbf{i den ordning entiteten deklarerar dem}. \code{display} deklarerar
    \code{number} och sedan \code{hex}, så det första objektet driver \code{number} och det andra
    drivs av \code{hex}.
  \item \code{input(7 downto 4)} är en \textbf{slice}, som instansen ser som en
    \code{std_logic_vector(3 downto 0)}.
\end{itemize}

Lägg märke till vad \code{hex_display}s arkitektur \emph{inte} innehåller: någon logik alls. Den
är ledningsdragning och inget annat. Varje grind i den färdiga designen kommer från de två
\code{display}-instanserna.

\subsection*{Positionell och namngiven associering}
\code{port map} ovan är \textbf{positionell} och matchar portar efter ordning. VHDL tillåter också
\textbf{namngiven} associering, som du kommer att möta i andras kod:

\begin{vhdlcode}
display1: entity work.display
    port map(number => input(7 downto 4), hex => hex1);
\end{vhdlcode}

De två är likvärdiga. Boken använder positionell överallt, vilket betyder att ordningen på en
entitets portdeklarationer är något du måste få rätt. Det är ingen formalitet: det är precis vad
de utdelade testbänkarna förlitar sig på, vilket är varför varje övnings självkontrollsruta
stavar ut sina portar \emph{i ordning}.

\subsection*{Vad det här kostar i hårdvara}
Inget överraskande. Att instansiera en modul två gånger lägger två kopior av dess grindar på
FPGA:n, precis som att skriva logiken två gånger skulle göra. Besparingen ligger i källkoden, inte
i kislet: en beskrivning, ett ställe att rätta på, och en design du kan läsa.
